\section{AIME I 2007}
        \begin{enumerate}
        	\item (\textit{AIME I 2007}) How many positive perfect squares less than $10^6$ are multiples of 24?


 
 \textbf{Answer: }083



\textbf{Solution}

The \href{https://artofproblemsolving.com/wiki/index.php?title=Prime\_factorization}{prime factorization} of $24$ is $2^3\cdot3$. Thus, each square must have at least $3$ factors of $2$ and $1$ factor of $3$ and its square root must have $2$ factors of $2$ and $1$ factor of $3$.This means that each square is in the form $(12c)^2$, where $12 c$ is a positive integer less than $\sqrt{10^6}$. There are $\left\lfloor \frac{1000}{12}\right\rfloor = \boxed{083}$ solutions.

 





\textbf{Solution 2}

The perfect squares divisible by $24$ are all multiples of $12$: $12^2$, $24^2$, $36^2$, $48^2$, etc... Since they all have to be less than $10^6$, or $1000^2$, the closest multiple of $12$ to $1000$ is $996$ ($12*83$), so we know that this is the last term in the sequence. Therefore, we know that there are $\boxed{083}$ perfect squares divisible by $24$ that are less than $10^6$.

 




	\item (\textit{AIME I 2007}) A 100 foot long moving walkway moves at a constant rate of 6 feet per second.  Al steps onto the start of the walkway and stands.  Bob steps onto the start of the walkway two seconds later and strolls forward along the walkway at a constant rate of 4 feet per second.  Two seconds after that, Cy reaches the start of the walkway and walks briskly forward beside the walkway at a constant rate of 8 feet per second.  At a certain time, one of these three persons is exactly halfway between the other two.  At that time, find the distance in feet between the start of the walkway and the middle person.


 
 \textbf{Answer: }052



\textbf{Solution}

Clearly we have people \textit{moving} at speeds of $6,8$ and $10$ feet/second. Notice that out of the three people, Cy is at the largest disadvantage to begin with and since all speeds are close, it is hardest for him to catch up. Furthermore, Bob is clearly the farthest along. Thus it is reasonable to assume that there is some point when Al is halfway between Cy and Bob. At this time $s$, we have that

 $\frac{8(s-4)+10(s-2)}{2}=6s$After solving, $s=\frac{26}{3}$. At this time, Al has traveled $6\cdot\frac{26}{3}=52$ feet.

 We could easily check that Al is in the middle by trying all three possible cases. $\frac{6s + 8(s-4)}{2} = 10(s-2)$ yields that $s = \frac 43$, which can be disregarded since both Bob and Cy hadn't started yet. $\frac{6s + 10(s-2)}{2} = 8(s-4)$ yields that $-10=-32$, a contradiction. Thus, the answer is $\boxed{052}$.

 




	\item (\textit{AIME I 2007}) The complex number $z$ is equal to $9+bi$, where $b$ is a positive real number and $i^{2}=-1$.  Given that the imaginary parts of $z^{2}$ and $z^{3}$ are the same, what is $b$ equal to?


 
 \textbf{Answer: }015



\textbf{Solution}

Squaring, we find that $(9 + bi)^2 = 81 + 18bi - b^2$. Cubing and ignoring the real parts of the result, we find that $(81 + 18bi - b^2)(9 + bi) = \ldots + (9\cdot 18 + 81)bi - b^3i$.

 Setting these two equal, we get that $18bi = 243bi - b^3i$, so $b(b^2 - 225) = 0$ and $b = -15, 0, 15$. Since $b > 0$, the solution is $\boxed{015}$.

 




	\item (\textit{AIME I 2007}) Three planets orbit a star circularly in the same plane.  Each moves in the same direction and moves at constant speed.  Their periods are $60$, $84$, and $140$ years.  The three planets and the star are currently collinear.  What is the fewest number of years from now that they will all be collinear again?


 
 \textbf{Answer: }105



\textbf{Solution}

Denote the planets $A, B, C$ respectively.  Let $a(t), b(t), c(t)$ denote the angle which each of the respective planets makes with its initial position after $t$ years.  These are given by $a(t) = \frac{t \pi}{30}$, $b(t) = \frac{t \pi}{42}$, $c(t) = \frac{t \pi}{70}$.

 In order for the planets and the central star to be collinear, $a(t)$, $b(t)$, and $c(t)$ must differ by a multiple of $\pi$.  Note that $a(t) - b(t) = \frac{t \pi}{105}$ and $b(t) - c(t) = \frac{t \pi}{105}$, so $a(t) - c(t) = \frac{ 2 t \pi}{105}$.  These are simultaneously multiples of $\pi$ exactly when $t$ is a multiple of $105$, so the planets and the star will next be collinear in $\boxed{105}$ years.

 




	\item (\textit{AIME I 2007}) The formula for converting a Fahrenheit temperature $F$ to the corresponding Celsius temperature $C$ is $C = \frac{5}{9}(F-32).$  An integer Fahrenheit temperature is converted to Celsius, rounded to the nearest integer, converted back to Fahrenheit, and again rounded to the nearest integer.


 For how many integer Fahrenheit temperatures between $32$ and $1000$ inclusive does the original temperature equal the final temperature?


 
 \textbf{Answer: }539



\textbf{Solution 1}

Examine $F - 32$ modulo 9.

 
                \begin{itemize}
                    \item If $F - 32 \equiv 0 \pmod{9}$, then we can define $9x = F - 32$. This shows that $F = \left[\frac{9}{5}\left[\frac{5}{9}(F-32)\right] + 32\right] \Longrightarrow F = \left[\frac{9}{5}(5x) + 32\right] \Longrightarrow F = 9x + 32$. This case works.\item If $F - 32 \equiv 1 \pmod{9}$, then we can define $9x + 1 = F - 32$. This shows that $F = \left[\frac{9}{5}\left[\frac{5}{9}(F-32)\right] + 32\right] \Longrightarrow F = \left[\frac{9}{5}(5x + 1) + 32\right] \Longrightarrow$$F = \left[9x + \frac{9}{5}+ 32 \right] \Longrightarrow F = 9x + 34$. So this case doesn't work.
                \end{itemize}


                Generalizing this, we define that $9x + k = F - 32$. Thus, $F = \left[\frac{9}{5}\left[\frac{5}{9}(9x + k)\right] + 32\right] \Longrightarrow F = \left[\frac{9}{5}(5x + \left[\frac{5}{9}k\right]) + 32\right] \Longrightarrow F = \left[\frac{9}{5} \left[\frac{5}{9}k \right] \right] + 9x + 32$. We need to find all values $0 \le k \le 8$ that $\left[ \frac{9}{5} \left[ \frac{5}{9} k \right] \right] = k$. Testing every value of $k$ shows that $k = 0, 2, 4, 5, 7$, so $5$ of every $9$ values of $k$ work.

 There are $\lfloor \frac{1000 - 32}{9} \rfloor = 107$ cycles of $9$, giving $5 \cdot 107 = 535$ numbers that work. Of the remaining $6$ numbers from $995$ onwards, $995,\ 997,\ 999,\ 1000$ work, giving us $535 + 4 = \boxed{539}$ as the solution.

 





\textbf{Solution 2}

Notice that $\left[ \frac{9}{5} \left[ \frac{5}{9} k \right] \right] = k$ holds if $k=\left[ \frac{9}{5}x\right]$ for some integer $x$.Thus, after translating from $F\to F-32$ we want count how many values of $x$ there are such that $k=\left[ \frac{9}{5}x\right]$ is an integer from $0$ to $968$. This value is computed as $\left[968*\frac{5}{9}\right]+1 = \boxed{539}$, adding in the extra solution corresponding to $0$.

 \textbf{Note} Proof that $\left[ \frac{9}{5} \left[ \frac{5}{9} k \right] \right] = k$ if $k=\left[ \frac{9}{5}x\right]$ for some integer $x$:

 First assume that $k$ cannot be written in the form $k=\left[ \frac{9}{5}x\right]$ for any integer $x$. Let $z = \left[ \frac{5}{9}k\right]$. Our equation simplifies to $k = \left[ \frac{9}{5}z\right]$. However, this equation is not possible, as we defined $k$ such that it could not be written in this form. Therefore, if $k \neq \left[ \frac{9}{5}x\right]$, then $\left[ \frac{9}{5} \left[ \frac{5}{9} k \right] \right] \neq k$.

 Now we will prove that if $k = \left[ \frac{9}{5}x\right]$, $\left[ \frac{9}{5} \left[ \frac{5}{9} k \right] \right] = k$. We realize that because of the 5 in the denominator of $\left[ \frac{9}{5}x \right]$, $\left[ \frac{9}{5}x \right]$ will be at most $\frac{2}{5}$ away from $\frac{9}{5}x$. Let $z = \left[ \frac{9}{5}x \right]- \frac{9}{5}x$, meaning that $-\frac{2}{5} \leq z \leq \frac{2}{5}$. Now we substitute this into our equation:

 
 \[\left[ \frac{9}{5} \left[ \frac{5}{9} k \right] \right] = \left[ \frac{9}{5} \left[ \frac{5}{9} \left[ \frac{9}{5}x\right] \right] \right] = \left[ \frac{9}{5} \left[ \frac{5}{9} (\frac{9}{5}x + z) \right] \right] = \left[ \frac{9}{5} \left[ \frac{5}{9} (\frac{9}{5}x + z) \right] \right] = \left[ \frac{9}{5} \left[ x+ \frac{5}{9}z \right] \right]\].

 Now we use the fact that $-\frac{2}{5} \leq z \leq \frac{2}{5}$

 
 \[\left[ \frac{9}{5} \left[ x - \frac{5}{9}(\frac{2}{5}) \right] \right] \leq \left[ \frac{9}{5} \left[ x + \frac{5}{9}(z) \right] \right] \leq \left[ \frac{9}{5} \left[ x + \frac{5}{9}(\frac{2}{5}) \right] \right]\]

 
 \[\left[ \frac{9}{5} x \right] \leq \left[ \frac{9}{5} \left[ x + \frac{5}{9}(z) \right] \right] = \left[ \frac{9}{5} \left[ \frac{5}{9}k \right] \right] \leq \left[ \frac{9}{5}x \right]\]

 Hence $\left[ \frac{9}{5} \left[ \frac{5}{9}k \right] \right] = \left[ \frac{9}{5}x \right] = k$ and we are done.

 - mako17

 





\textbf{Solution 3}

Let $c$ be a degree Celsius, and $f=\frac 95c+32$ rounded to the nearest integer. Since $f$ was rounded to the nearest integer we have $|f-((\frac 95)c+32)|\leq 1/2$, which is equivalent to $|(\frac 59)(f-32)-c|\leq \frac 5{18}$ if we multiply by $5/9$. Therefore, it must round to $c$ because $\frac 5{18}<\frac 12$ so $c$ is the closest integer. Therefore there is one solution per degree Celsius in the range from $0$ to $(\frac 59)(1000-32) + 1=(\frac 59)(968) + 1=538.8$, meaning there are $539$ solutions.

 





\textbf{Solution 4}

Start listing out values for $F$ and their corresponding values of $C$. You will soon find that every 9 values starting from $F$ = 32, there is a pattern:

 $F=32$: Works

 $F=33$: Doesn't work

 $F=34$: work

 $F=35$: Doesn’t work

 $F=36$: Works

 $F=37$: Works

 $F=38$: Doesn’t work

 $F=39$: Works

 $F=40$: Doesn’t work

 $F=41$: Works

 There are $969$ numbers between $32$ and $1000$, inclusive. This is $107$ sets of $9$, plus $6$ extra numbers at the end. In each set of $9$, there are $5$ “Works,” so we have $107\cdot5 = 535$ values of $F$ that work.

 Now we must add the $6$ extra numbers. The number of “Works” in the first $6$ terms of the pattern is $4$, so our final answer is $535 + 4 = 539$ solutions that work.

 Submitted by warriorcats

 





\textbf{Solution 5(similar to solution 3 but faster solution if you have no time)}

Notice that every $C$ value corresponds to exactly one $F$ value but multiple $F$ values can correspond to a $C$ value. Thus, the smallest $C$ value is $0$ and the largest $C$ value is $538$ yielding $\boxed{539}$ solutions.

 -alanisawesome2018

 \href{https://mathproblemsolvingskills.wordpress.com/}{MathProblemSolvingSkills.com}

 


 




	\item (\textit{AIME I 2007}) A frog is placed at the \href{https://artofproblemsolving.com/wiki/index.php?title=Origin}{origin} on the \href{https://artofproblemsolving.com/wiki/index.php?title=Number\_line}{number line}, and moves according to the following rule: in a given move, the frog advances to either the closest \href{https://artofproblemsolving.com/wiki/index.php?title=Point}{point} with a greater \href{https://artofproblemsolving.com/wiki/index.php?title=Integer}{integer} \href{https://artofproblemsolving.com/wiki/index.php?title=Coordinate}{coordinate} that is a multiple of $3$, or to the closest point with a greater integer coordinate that is a multiple of $13$. A \textit{move sequence} is a \href{https://artofproblemsolving.com/wiki/index.php?title=Sequence}{sequence} of coordinates which correspond to valid moves, beginning with $0$, and ending with $39$. For example, $0,\ 3,\ 6,\ 13,\ 15,\ 26,\ 39$ is a move sequence. How many move sequences are possible for the frog?


 
 \textbf{Answer: }169



\textbf{Solution 1}

Let us keep a careful tree of the possible number of paths around every multiple of $13$.

 From $0 \Rightarrow 13$, we can end at either $12$ (mult. of 3) or $13$ (mult. of 13). 

 
                \begin{itemize}
                    <li>Only $1$ path leads to $12$
                \begin{itemize}
                    \item Continuing from $12$, there is $1 \cdot 1 = 1$ way to continue to $24$\item There are $1 \cdot \left(\frac{24-15}{3} + 1\right) = 4$ ways to reach $26$.
                \end{itemize}


                </li><li>There are $\frac{12 - 0}{3} + 1 = 5$ ways to reach $13$.
                \begin{itemize}
                    \item Continuing from $13$, there are $5 \cdot 1 = 5$ ways to get to $24$\item There are $5 \cdot \left(\frac{24-15}{3} + 1 + 1\right) = 25$ ways (the first 1 to make it inclusive, the second to also jump from $13 \Rightarrow 26$) to get to $26$.
                \end{itemize}


                </li>
                \end{itemize}


                Regrouping, work from $24 | 26\Rightarrow 39$

 
                \begin{itemize}
                    <li>There are $1 + 5 = 6$ ways to get to $24$
                \begin{itemize}
                    \item Continuing from $24$, there are $6 \cdot \left(\frac{39 - 27}{3}\right) = 24$ ways to continue to $39$.
                \end{itemize}


                </li><li>There are $4 + 25 = 29$ ways to reach $26$.
                \begin{itemize}
                    \item Continuing from $26$, there are $29 \cdot \left(\frac{39-27}{3} + 1\right) = 145$ (note that the 1 is not to inclusive, but to count $26 \Rightarrow 39$).
                \end{itemize}


                </li>
                \end{itemize}


                In total, we get $145 + 24 = 169$.

 


 In summary, we can draw the following tree, where in $(x,y)$, $x$ represents the current position on the number line, and $y$ represents the number of paths to get there:

 




	\item (\textit{AIME I 2007}) Let 
$N = \sum_{k = 1}^{1000} k ( \lceil \log_{\sqrt{2}} k \rceil  - \lfloor \log_{\sqrt{2}} k \rfloor )$


 Find the remainder when $N$ is divided by 1000.  ($\lfloor{k}\rfloor$ is the greatest integer less than or equal to $k$, and $\lceil{k}\rceil$ is the least integer greater than or equal to $k$.)


 
 \textbf{Answer: }477



\textbf{Solution}

The ceiling of a number minus the floor of a number is either equal to zero (if the number is an \href{https://artofproblemsolving.com/wiki/index.php?title=Integer}{integer}); otherwise, it is equal to 1. Thus, we need to find when or not $\log_{\sqrt{2}} k$ is an integer.

 The change of base formula shows that $\frac{\log k}{\log \sqrt{2}} = \frac{2 \log k}{\log 2}$. For the $\log 2$ term to cancel out, $k$ is a \href{https://artofproblemsolving.com/wiki/index.php?title=Exponent}{power} of $2$. Thus, $N$ is equal to the sum of all the numbers from 1 to 1000, excluding all powers of 2 from $2^0 = 1$ to $2^9 = 512$. 

 The formula for the sum of an \href{https://artofproblemsolving.com/wiki/index.php?title=Arithmetic\_sequence}{arithmetic sequence} and the sum of a \href{https://artofproblemsolving.com/wiki/index.php?title=Geometric\_sequence}{geometric sequence} yields that our answer is $\left[\frac{(1000 + 1)(1000)}{2} - (1 + 2 + 2^2 + \ldots + 2^9)\right] \mod{1000}$.

 Simplifying, we get$\left[1000\left(\frac{1000+1}{2}\right) -1023\right] \mod{1000} \equiv [500-23] \mod{1000} \equiv 477 \mod{1000}.$ The answer is $\boxed{477}$

 




	\item (\textit{AIME I 2007}) The polynomial $P(x)$ is cubic.  What is the largest value of $k$ for which the polynomials $Q_1(x) = x^2 + (k-29)x - k$ and $Q_2(x) = 2x^2+ (2k-43)x + k$ are both factors of $P(x)$?


 
 \textbf{Answer: }030



\textbf{Solution 1}

We can see that $Q_1$ and $Q_2$ must have a \href{https://artofproblemsolving.com/wiki/index.php?title=Root}{root} in common for them to both be \href{https://artofproblemsolving.com/wiki/index.php?title=Factor}{factors} of the same cubic.

 Let this root be $a$.

 We then know that $a$ is a root of$Q_{2}(x)-2Q_{1}(x) = 2x^{2}+2kx-43x+k-2x^{2}-2kx+58x+2k = 15x+3k$, so $Q_{2}(a)-2Q_1(a)=15a+3k=0\implies a = \frac{-k}{5}$.

 We then know that $\frac{-k}{5}$ is a root of $Q_{1}$ so we get:$\frac{k^{2}}{25}+(k-29)\left(\frac{-k}{5}\right)-k = 0 = k^{2}-5(k-29)(k)-25k = k^{2}-5k^{2}+145k-25k$or $k^{2}=30k$, so $k=30$ is the highest.

 We can trivially check into the original equations to find that $k=30$ produces a root in common, so the answer is $\boxed{030}$.

 \textbf{Proof of the common root} Factor $P(x)$ as $(x-a)(x-b)(x-c)$. Then, for $Q_1(x)$ and $Q_2(x)$, both are quadratics, to properly divide into $P(x)$, both roots must be from $P(x)$ to properly cancel. Therefore, 2 possible roots exist for $Q_1(x)$ and 2 for $Q_2(x)$, 4 roots total to be chosen from the 3 plausible roots. That is the same as putting 4 balls into 3 bins. Through the \href{https://artofproblemsolving.com/wiki/index.php?title=Pigeonhole\_Principle}{Pigeonhole Principle}, we have that $\lfloor \frac{4-1}{1} \rfloor + 1 = 1+1 = 2$ balls will end up in one bin, meaning a root must be shared. 

 We now prove this must be the case. If the two balls were from the same quadratic, $P$ would have to be $(x-a)^2(x-b)(x-c)$ which is a quartic and impossible, as each $Q_i$ must fully divide the three roots of $P$, so we are done. $\square$

 





\textbf{Solution 2}

Again, let the common root be $a$; let the other two roots be $m$ and $n$. We can write that $(x - a)(x - m) = x^2 + (k - 29)x - k$ and that $2(x - a)(x - n) = 2\left(x^2 + \left(k - \frac{43}{2}\right)x + \frac{k}{2}\right)$.

 Therefore, we can write four equations (and we have four \href{https://artofproblemsolving.com/wiki/index.php?title=Variable}{variables}), $a + m = 29 - k$, $a + n = \frac{43}{2} - k$, $am = -k$, and $an = \frac{k}{2}$. 

 The first two equations show that $m - n = 29 - \frac{43}{2} = \frac{15}{2}$. The last two equations show that $\frac{m}{n} = -2$. Solving these show that $m = 5$ and that $n = -\frac{5}{2}$. Substituting back into the equations, we eventually find that $k = \boxed{030}$.

 





\textbf{Solution 3}

Since $Q_1(x)$ and $Q_2(x)$ are both factors of $P(x)$, which is cubic, we know the other factors associated with each of $Q_1(x)$ and $Q_2(x)$ must be linear. Let $Q_1(x)R(x) = Q_2(x)S(x) = P(x)$, where $R(x) = ax + b$ and $S(x) = cx + d$. Then we have that $((x^2 + (k-29)x - k))(ax + b) = ((2x^2+ (2k-43)x + k))(cx + d)$. Equating coefficients, we get the following system of equations: 

 
 \begin{align}  a = 2c \\ b = -d \\ 2c(k - 29) - d = c(2k - 43) + 2d \\ -d(k - 29) - 2ck = d(2k - 43) + ck  \end{align}

 Using equations $(1)$ and $(2)$ to make substitutions into equation $(3)$, we see that the $k$'s drop out and we're left with $d = -5c$. Substituting this expression for $d$ into equation $(4)$ and solving, we see that $k$ must be $\boxed {030}$.  

 





\textbf{Solution 4}

Notice that if the roots of $Q_1(x)$ and $Q_2(x)$ are all distinct, then $P(x)$ would have four distinct roots, which is a contradiction since it's cubic. 

 Thus, $Q_1(x)$ and $Q_2(x)$ must share a root. Let this common value be $r.$ 

 Thus, we see that we have 
 \[r^2 + (k - 29)r - k = 0,\]
 \[2r^2 + (2k - 43)r + k = 0.\] Adding the two equations gives us 
 \[3r^2 + (3k - 72)r = 0 \implies r = 0, 24 - k.\] Now, we have two cases to consider. If $r = 0,$ then we have that $Q_1(r) = 0 = r^2 + (k - 29)r - k \implies k = 0.$ On the other hand, if $r = 24 - k,$ we see that 
 \[Q_1(r) = 0 = (24 - k)^2 + (k - 29)(24 - k) - k \implies k = \boxed{030}.\] This can easily be checked to see that it does indeed work, and we're done! 

 





\textbf{Solution 5}

Since $Q_1(x) = x^2 + (k - 29)x - k$ and $Q_2(x) = 2x^2 + (2k - 43)x + k$ are both factors of the cubic polynomial $P(x)$, and both are quadratics, the only way for them to divide the same cubic is if they share a common linear factor. Otherwise, their product would be degree 4, which is too large. Therefore, $P(x)$ can be factored as$P(x) = Q_1(x)(2x + a) = Q_2(x)(x + b)$for some constants $a, b$. We choose the linear factors so that both products have degree 3 with leading coefficient 2: since $Q_1$ has leading coefficient 1, its linear factor must have leading coefficient 2, and since $Q_2$ has leading coefficient 2, its linear factor must have leading coefficient 1.

 Looking at constant terms on both sides,$-k \cdot a = k \cdot b$,so$-k a = k b \implies a = -b.$

 Next, expanding both sides and comparing the $x^2$ coefficients gives$a + 2k - 58 = 2b + 2k - 43.$Substituting $a = -b$, we get$-b + 2k - 58 = 2b + 2k - 43 \implies -b = 2b + 15 \implies b = -5, \quad a = 5.$

 Finally, comparing the $x$ coefficients:$a(k - 29) - 2k = (2k - 43) b + k,$which becomes$5(k - 29) - 2k = (2k - 43)(-5) + k,$simplifying to$3k - 145 = -9k + 215.$Solving for $k$ yields$12k = 360 \implies \boxed{30}.$

 




	\item (\textit{AIME I 2007}) In right triangle $ABC$ with right angle $C$, $CA = 30$ and $CB = 16$.  Its legs $CA$ and $CB$ are extended beyond $A$ and $B$.  Points $O_1$ and $O_2$ lie in the exterior of the triangle and are the centers of two circles with equal radii.  The circle with center $O_1$ is tangent to the hypotenuse and to the extension of leg $CA$, the circle with center $O_2$ is tangent to the hypotenuse and to the extension of leg $CB$, and the circles are externally tangent to each other.  The length of the radius of either circle can be expressed as $p/q$, where $p$ and $q$ are relatively prime positive integers.  Find $p+q$.


 
 \textbf{Answer: }737



\textbf{Solution 1}


                    \begin{figure}[H]
                        \centering
                        \includegraphics[width=0.5\textwidth]{Solution/AIME_2007_I_9_Solution_Figure_1.png}
                    \end{figure}
                    

 Label the points as in the diagram above. If we draw $\overline{O_1A}$ and $\overline{O_2B}$, we form two \href{https://artofproblemsolving.com/wiki/index.php?title=Right\_triangle}{right triangles}. As $\overline{AF}$ and $\overline{AD}$ are both \href{https://artofproblemsolving.com/wiki/index.php?title=Tangent}{tangents} to the circle, we see that $\overline{O_1A}$ is an \href{https://artofproblemsolving.com/wiki/index.php?title=Angle\_bisector}{angle bisector}. Thus, $\triangle AFO_1 \cong \triangle ADO_1$. Call $x = AD = AF$ and $y = EB = BG$. We know that $x + y + 2r = 34$. 

 If we call $\angle CAB = \theta$, then $\angle DAO_1 = \frac{180 - \theta}{2}$. Apply the \href{https://artofproblemsolving.com/wiki/index.php?title=Trigonometric\_identities}{tangent half-angle formula} ($\tan \frac{\theta}{2} = \sqrt{\frac{1 - \cos \theta}{1 + \cos \theta}}$). We see that $\frac rx = \tan \frac{180 - \theta}{2} = \sqrt{\frac{1 - \cos (180 - \theta)}{1 + \cos (180 - \theta)}}$$= \sqrt{\frac{1 + \cos \theta}{1 - \cos \theta}}$. Also, $\cos \theta = \frac{30}{34} = \frac{15}{17}$. Thus, $\frac rx = \sqrt{\frac{1 + \frac{15}{17}}{1 - \frac{15}{17}}}$, and $x = \frac{r}{4}$. 

 Similarly, we find that $y = r/\sqrt{\frac{1 + \frac{8}{17}}{1 - \frac{8}{17}}} = \frac{3r}{5}$.

 Therefore, $x + y + 2r = \frac{r}{4} + \frac{3r}{5} + 2r = \frac{57r}{20} = 34 \Longrightarrow r = \frac{680}{57}$, and $p + q = 737$.

 





\textbf{Solution 2}

Use a similar solution to the aforementioned solution. Instead, call $\angle CAB = 2\theta$, and then proceed by simplifying through identities. We see that $\frac rx = \tan \left(\frac{180 - 2\theta}{2}\right) = \tan (90 - \theta)$. In terms of $r$, we find that $x = \frac{r}{\cot \theta} = \frac{r\sin \theta}{\cos \theta}$. Similarly, we find that $y = \frac{r \sin(45 - \theta)}{\cos (45 - \theta)}$. 

 Substituting, we find that $r\left(\frac{\sin \theta}{\cos \theta} + \frac{\sin(45 - \theta)}{\cos (45 - \theta)} + 2\right) = 34$. Under a common denominator, $r\left(\frac{\sin \theta \cos (45 - \theta) + \cos \theta \sin (45 - \theta)}{\cos \theta \cos (45 - \theta)} + 2\right) = 34$. \href{https://artofproblemsolving.com/wiki/index.php?title=Trigonometric\_identities}{Trigonometric identities} simplify this to $r\left(\frac{\sin\left((\theta) + (45 - \theta)\right)}{\frac 12 \left(\cos (\theta + 45 - \theta) + \cos (\theta - 45 + \theta) \right)} + 2\right) = 34$. From here, it is possible to simplify:

 
            \begin{enumerate}[label=(\roman*)]
                \item $r\left(\frac{2 \sin 45}{\cos 45 + \cos 2\theta \cos 45 + \sin 2\theta \sin 45} +2\right) = 34$\item $r\left(\frac{2}{\frac{17}{17} + \frac{8}{17} + \frac{15}{17}} + 2\right) = 34$\item $r\left(\frac{57}{20}\right) = 34$
            \end{enumerate}
            The answer is $34 \cdot \frac{20}{57} = \frac{680}{57}$, and $p + q = 737$.

 





\textbf{Solution 3}

Let the point where CB's extension hits the circle be G, and the point where the hypotenuse hits that circle be E. Clearly $EB=GB$. Let $EB=x$. Draw the two \href{https://artofproblemsolving.com/wiki/index.php?title=Perpendicular}{perpendicular} radii to G and E. Now we have a \href{https://artofproblemsolving.com/wiki/index.php?title=Cyclic\_quadrilateral}{cyclic quadrilateral}. Let the radius be length $r$. We see that since the cosine of angle ABC is $\frac{15}{17}$ the cosine of angle EBG is $-\frac{15}{17}$. Since the measure of the angle opposite to EBG is the \href{https://artofproblemsolving.com/wiki/index.php?title=Complement}{complement} of this one, its cosine is $\frac{15}{17}$. Using the law of cosines, we see that $x^{2}+x^{2}+\frac{30x^{2}}{17}=r^{2}+r^{2}-\frac{30r^{2}}{17}$This tells us that $r=4x$.

 Now look at the other end of the hypotenuse. Call the point where CA hits the circle F and the point where the hypotenuse hits the circle D. Draw the radii to F and D and we have cyclic quadrilaterals once more.Using the law of cosines again, we find that the length of our tangents is $2.4x$. Note that if we connect the centers of the circles we have a rectangle with sidelengths 8x and 4x. So, $8x+2.4x+x=34$. Solving we find that $4x=\frac{680}{57}$ so our answer is 737.

 





\textbf{Solution 4}


                    \begin{figure}[H]
                        \centering
                        \includegraphics[width=0.5\textwidth]{Solution/AIME_2007_I_9_Solution_Figure_1.png}
                    \end{figure}
                    By Pythagoras, $AB = 34$.  Let $I_{C}$ be the $C$-excenter of triangle $ABC$.  Then the $C$-exradius $r_{C}$ is given by$r_{C}= \frac{K}{s-c}= \frac{240}{40-34}= 40$.

 The circle with center $O_{1}$ is tangent to both $AB$ and $AC$, which means that $O_{1}$ lies on the external angle bisector of $\angle BAC$.  Therefore, $O_{1}$ lies on $AI_{C}$.  Similarly, $O_{2}$ lies on $BI_{C}$.

 Let $r$ be the common radius of the circles with centers $O_{1}$ and $O_{2}$.  The distances from points $O_{1}$ and $O_{2}$ to $AB$ are both $r$, so $O_{1}O_{2}$ is parallel to $AB$, which means that triangles $I_{C}AB$ and $I_{C}O_{1}O_{2}$ are similar.

 The distance from $I_{C}$ to $AB$ is $r_{C}= 40$, so the distance from $I_{C}$ to $O_{1}O_{2}$ is $40-r$.  Therefore,

 $\frac{40-r}{40}= \frac{O_{1}O_{2}}{AB}= \frac{2r}{34}\quad \Rightarrow \quad r = \frac{680}{57}$.

 Hence, the final answer is $680+57 = 737$.

 





\textbf{Solution 5}


                    \begin{figure}[H]
                        \centering
                        \includegraphics[width=0.5\textwidth]{Solution/AIME_2007_I_9_Solution_Figure_1.png}
                    \end{figure}
                    Start with a scaled 16-30-34 triangle. \href{https://artofproblemsolving.com/wiki/index.php?title=Inscribe}{Inscribe} a circle.  The height, $h,$ and radius, $r,$ are found via $A=\frac{1}{2}\times 16s\times 30s=\frac{1}{2}\times 34s\times h=\frac{1}{2}\times rp,$ where $p$ is the \href{https://artofproblemsolving.com/wiki/index.php?title=Perimeter}{perimeter}.

 Cut the figure through the circle and perpendicular to the hypotenuse.  Slide the two pieces in opposite directions along the hypotenuse until they are one diameter of the circle apart.  Complete the two partial circles.

 The linear dimensions of the new triangle are $\frac{46s}{34s}=\frac{23}{17}$ times the size of the original.  The problem's 16-30-34 triangle sits above the circles.  Equate heights and solve for $r=6s$:

 $\frac{240s}{17}\times\frac{23}{17} =  \frac{240}{17}+12s$

 $20s\times 23 =  20\times 17+s\times 17\times 17$

 $s = \frac{340}{171}$

 $r = 6s = \frac{680}{57}$The answer is $737$.

 





\textbf{Solution 6}


                    \begin{figure}[H]
                        \centering
                        \includegraphics[width=0.5\textwidth]{Solution/AIME_2007_I_9_Solution_Figure_1.png}
                    \end{figure}
                    

 Using \href{https://artofproblemsolving.com/wiki/index.php?title=Homothety}{homothety} in the diagram above, as well as the \href{https://artofproblemsolving.com/wiki/index.php?title=Auxiliary&action=edit&redlink=1}{auxiliary} triangle, leads to the solution.

 





\textbf{Solution 7}

A different approach is to plot the triangle on the Cartesian Plane with $C$ at $(0,0)$, $A$ at $(0,30)$, and $B$ at $(16,0)$. We wish to find the coordinates of $O_1$ and $O_2$ in terms of the radius, which will be expressed as $r$ in the rest of this solution. When we know the coordinates, we will set the distance between the 2 points equal to $2r$. All points $r$ units away from $\overline{AB}$ are on the line with slope $-\frac{15}{8}$, and y-intercept $30+ \frac{17}{8} r$

 $O_1$ will have x-coordinate $r$ and likewise $O_2$ will have y-coordinate $r$plugging this into the equation for the line mentioned in the sentence above gives us:

 $O_1 = \left(r,\frac14 r+30\right)$ and $O_2 = \left(\frac35 r+16,r\right)$

 By the distance formula and the fact that the circles and tangent, we have: $\left(16-\frac25 r\right)^2 + \left(30-\frac34 r\right)^2 = (2r)^2$

 which simplifies into the quadratic equation: $1311 r^2 + 23120 r - 462400 = 0$

 And by the quadratic equation, the solutions are: $\frac{-23120 \pm 54400}{2622}$The solution including the "$-$" is extraneous so we have the radius equal to$\frac{31280}{2622}$

 Which simplifies to $\frac{680}{57}$. The sum of the numerator and the denominator is $\boxed{737}$

 





\textbf{Solution 8 (simple algebra)}

It is known that $O_1O_2$ is parallel to AB. Thus, extending $O_1F$ and $GO_2$ to intersect at H yields similar triangles $O_1O_2H$ and BAC, so that $O_1O_2 = 2r$, $O_1H = \frac{16r}{17}$, and $HO_2 = \frac{30r}{17}$. It should be noted that $O_2G = r$. Also, FHGC is a rectangle, and so AF = $\frac{47r}{17} - 30$ and similarly for BG. Because tangents to a circle are equal, the hypotenuse can be expressed in terms of r: 
 \[2r + \frac{47r}{17} - 30 + \frac{33r}{17} - 16 = 34\] Thus, r = $\frac{680}{57}$, and the answer is $\boxed{737}.$

 \textbf{Note} When drawing the diagram, it may seem that $H$ lies on circle $O_1$, but it is actually not: $H$ lies inside of circle $O_1$. We can see this from the similarity ratios: $\frac{O_1O_2}{BA}=\frac{HO_1}{CB}=\frac{HO_2}{CA}$. Taking a look at the first equation ($\frac{O_1O_2}{BA}=\frac{HO_1}{CB}$), $\frac{2r}{34}=\frac{HO_1}{16}$ which simplifies to $\frac r{17}=\frac{HO_1}{16}$. Indeed, $HO_1$ does not equal $r$, instead, $HO_1=\frac{16}{17}r$. 

 





\textbf{Solution 9}

Let the radius of the circle be $r$. It can be seen that $\Delta FHO_{1}$ and $\Delta O_{2}GJ$ are similar to $\Delta ACB$, and the length of the hypotenuses are $\frac{17}{8}r$ and $\frac {17}{15}r$, respectively. Then, the entire length of $HJ$ is going to be $(\frac{17}{8}+\frac{17}{15}+2)r = \frac{631}{120}r$. The length of the hypotenuse of $\Delta ACB$ is 34, so the length of the height to $AB$ is $\frac{16*30}{34} = \frac{240}{17}$. Thus, the height to $\Delta HCJ$ is going to be $\frac{240}{17} + r$. $\Delta HCJ$ is similar to $\Delta ACB$ so we have the following: $\frac{\frac{631}{120}r}{34} = \frac{\frac{240}{17} + r}{\frac{240}{17}}$. Cross multiplying and simplifying, we get that $r = \frac{680}{57}$ so the answer is $\boxed{737}$.~Leonard\_my\_dude~

 


 \href{https://mathproblemsolvingskills.wordpress.com/}{MathProblemSolvingSkills.com}

 


 




	\item (\textit{AIME I 2007}) In a $6 \times 4$ grid ($6$ rows, $4$ columns), $12$ of the $24$ squares are to be shaded so that there are two shaded squares in each row and three shaded squares in each column.  Let $N$ be the number of shadings with this property.  Find the remainder when $N$ is divided by $1000$.


 
                    \begin{figure}[H]
                        \centering
                        \includegraphics[width=0.5\textwidth]{Problem/AIME_2007_I_10_Problem_Figure_1.png}
                    \end{figure}
                    


 
 \textbf{Answer: }860



\textbf{Solution 1}

Consider the first column.  There are ${6\choose3} = 20$ ways that the rows could be chosen, but without loss of generality let them be the first three rows.  (Change the order of the rows to make this true.)  We will multiply whatever answer we get by 20 to get our final answer.

 Now consider the 3x3 that is next to the 3 boxes we have filled in.  We must put one ball in each row (since there must be 2 balls in each row and we've already put one in each).  We split into three cases:

 
                \begin{itemize}
                    \item All three balls are in the same column.  In this case, there are $3$ choices for which column that is.  From here, the bottom half of the board is fixed.
                \end{itemize}


                
                \begin{itemize}
                    \item Two balls are in one column, and one is in the other.  In this case, there are 3 ways to choose which column gets 2 balls and 2 ways to choose which one gets the other ball.  Then, there are 3 ways to choose which row the lone ball is in.  Now, what happens in the bottom half of the board?  Well, the 3 boxes in the column with no balls in the top half must all be filled in, so there are no choices here.  In the column with two balls already, we can choose any of the 3 boxes for the third ball.  This forces the location for the last two balls.  So we have $3 \cdot 2 \cdot 3 \cdot 3 = 54$.
                \end{itemize}


                
                \begin{itemize}
                    \item All three balls are in different columns.  Then there are 3 ways to choose which row the ball in column 2 goes and 2 ways to choose where the ball in column 3 goes.  (The location of the ball in column 4 is forced.)  Again, we think about what happens in the bottom half of the board.  There is 1 ball in each of the columns 1, 2, and 3 now, so in the 3x3 where we still have choices, each row and column has one square that is \textbf{not} filled in. Similar to the upper 3x3, there are 6 ways to do this.  So there are $3 \cdot 2 \cdot 6 = 36$ ways.
                \end{itemize}


                Therefore there are $20(3+54+36) = 1860$ different shadings, and the solution is $\boxed{860}$.

 





\textbf{Solution 2}

We start by showing that every group of $6$ rows can be grouped into $3$ complementary pairs.We proceed with proof by contradiction. Without loss of generality, assume that the first row has columns $1$ and $2$ shaded. Note how if there is no complement to this, then all the other five rows must have at least one square in the first two columns shaded. That means that in total, the first two rows have $2+5=7$ squares shaded in- that is false since it should only be $6$.Thus, there exists another row that is complementary to the first. We remove those two and use a similar argument again to show that every group of $6$ rows can be grouped into $3$ complementary pairs.

 Now we proceed with three cases.

 
                \begin{itemize}
                    \item There are $\frac{\binom42}2=3$ pairs of complementary pairs. The first case is that the three pairs are all different, meaning that every single possible pair of shaded squares is used once. This gives us $6!=720.$
                \end{itemize}


                
                \begin{itemize}
                    \item Our second case is that two of the pairs are the same, and the third is different. We have $3$ to choose the pair that shows up twice and $2$ for the other, giving us $3\cdot2\cdot\binom62\binom42\binom21=6\cdot15\cdot6\cdot2=1080.$
                \end{itemize}


                
                \begin{itemize}
                    \item Our final case is that all three pairs are the same. This is just $3\cdot\binom63=60.$
                \end{itemize}


                Our answer is thus $720+1080+60=1860,$ leaving us with a final answer of $\boxed{860}.$

 





\textbf{Solution 3}

We draw a \href{https://artofproblemsolving.com/wiki/index.php?title=Bijection}{bijection} between walking from $(0,0,0,0)$ to $(3,3,3,3)$ as follows: if in the $i$th row, the $j$th and $k$th columns are shaded, then the $(2i-1)$st step is in the direction corresponding to $j$, and the $(2i)$th step is in the direction corresponding to $k$ ($j < k$) here. We can now use the \href{https://artofproblemsolving.com/wiki/index.php?title=Principle\_of\_Inclusion-Exclusion}{Principle of Inclusion-Exclusion} based on the stipulation that $j\ne k$ to solve the problem:

 $\frac{1}{2^6}\left(\frac{12!}{3!^4}-4\cdot{6 \choose 1}\frac{10!}{3!^3}+4\cdot 3\cdot{6 \choose 2}\frac{8!}{3!^2}-4\cdot 3\cdot 2\cdot{6 \choose 3}\frac{6!}{3!}+4\cdot 3\cdot 2\cdot 1\cdot {6 \choose 4}4!\right) = 1860$

 So that the answer is $860$.

 





\textbf{Solution 4}

There are ${6\choose3}$ to choose the arrangement of the shaded squares in each column. Examine the positioning of the shaded squares in the first two columns:

 
                    \begin{figure}[H]
                        \centering
                        \includegraphics[width=0.5\textwidth]{Solution/AIME_2007_I_10_Solution_Figure_1.PNG}
                    \end{figure}
                    One example of each case for the first two columns


 


 


 


 


 


 


 
                \begin{itemize}
                    \item If column 1 and column 2 do not share any two filled squares on the same row, then there are ${6\choose3}$ combinations for column 1, and then column 2 is fixed. Now, any row cannot have more than 2 shaded squares, so after we pick three more squares in the third column, the fourth column is also fixed. This gives ${6\choose3}^2 = 400$ arrangements.\item If column 1 and column 2 share 1 filled square on the same row (6 places), then they each share 1 filled square on a row ($6 - 1 = 5$ places), share another empty square on a row, and have 2 squares each on different rows. This gives $6 \cdot 5 \cdot {4\choose2} = 180$. Now, the third and fourths columns must also share a fixed shared shaded square in the row in which the first two columns both had spaces, and another fixed empty square. The remaining shaded squares can only go in 4 places, so we get ${4\choose2} = 6$. We get $180 \cdot 6 = 1080$.\item If column 1 and column 2 share 2 filled squares on the same row (${6\choose2} = 15$ places), they must also share 2 empty squares on the same row (${4\choose2} = 6$). The last two squares can be arranged in ${2\choose1} = 2$ positions; this totals to $15 \cdot 6 \cdot 2 = 180$. Now, the third and fourth columns have a fixed 2 filled squares in common rows and 2 empty squares in common rows. The remaining 2 squares have ${2\choose1} = 2$ places, giving $180 \cdot 2 = 360$.\item If column 1 and column 2 share 3 filled squares on the same row (${6\choose3} = 20$ places), then the squares on columns 3 and 4 are fixed.
                \end{itemize}


                Thus, there are $400 + 1080 + 360 + 20 = 1860$ number of shadings, so the answer is $860$.

 





\textbf{Solution 5}

Consider all possible shadings for a single row. There are ${4 \choose 2}=6$ ways to do so, and denote these as $a=1+2$, $b=3+4$, $c=1+4$, $d=2+3$, $e=1+3$, and $f=2+4$ where $x+y$ indicates that columns $x$ and $y$ are shaded. From our condition on the columns, we have $a+c+e=a+d+f=b+d+e=b+c+f=3$ Summing the first two and the last two equations, we have $2a+c+d+e+f=6=2b+c+d+e+f$, from which we have $a=b$. Likewise, $c=d$ and $e=f$ since these pairs shade in complimentary columns. So the six rows are paired up into a row and its compliment. In all, we can have 3 a's and 3 b's and similar setups for $c$/$d$ and $e$/$f$, 2 a's, 2 b's, 1 c and 1 d and similar setups for all six arrangements, or one of each. This first case gives $3{6 \choose 3}=60$ solutions; the second gives $6\cdot 6\cdot5\cdot{4 \choose 2}=1080$ solutions, and the final case gives $6!=720$ solutions. In all, we have 1860 solutions, for an answer of $860$.

 





\textbf{Solution 6}

Each shading can be brought, via row swapping operations, to a state with a $3\times2$ shaded $L$ in the lower left hand corner. The number of such arrangements multiplied by ${5 \choose 2}{3\choose 2}$ will be the total. Consider rows 2 and 3 up from the bottom: they each have one of their allotted two squares shaded. Depending how the remaining three shades are distributed, the column totals of columns 2,3, and 4 from the left can be of the form $\{3,0,0\},\{0,1,2\},\{1,1,1\}$.Form 1: The entire lower left $3\times2$ rectangle is shaded, forcing the opposite $3\times2$ rectangle to also be shaded; thus 1 arrangementForm 2: There is a column with nothing shaded in the bottom right $3\times2$, so it must be completely shaded in the upper right $3\times2$. Now consider the upper right half column that will have $1$ shade. There are $3$ ways of choosing this shade, and all else is determined from here; thus 3 arrangementsForm 3: The upper right $3\times3$ will have exactly $2$ shades per column and row. This is equivalent to the number of terms in a $3\times3$ determinant, or $6$ arrangements

 Of the $3^2$ ways of choosing to complete the bottom half of the $4\times6$, form 1 is achieved in exactly 1 way; form 2 is achieved in $3\times2$ ways; and form $3$ in the remaining $2$ ways.Thus, the weighted total is $1+6\times3+2\times6=31$.Complete: $31\times60=860\mod{1000}$

 





\textbf{Solution 7}

Note that if we find a valid shading of the first 3 columns, the shading of the last column is determined. We also note that within the first 3 columns, there will be 3 rows with 1 shaded square and 3 rows with 2 shaded squares.

 There are ${6 \choose 3}$ ways to choose which rows have 1 shaded square (which we'll call a "1-row") within the first 3 columns and which rows have 2 (we'll call these "2-rows") within the first 3 columns. Next, we do some casework:

 
                \begin{itemize}
                    \item If all of the shaded squares in the first column are in a 1-row, then the squares in the second and third columns must be in the 2-rows. Thus there is only ${3 \choose 3}{3 \choose 0} \times 1= 1$ valid shading in this case.
                \end{itemize}


                
                \begin{itemize}
                    \item If 2 of the shaded squares in the first column are in a 1-row, and the third shaded square is in a 2-row, then the other shaded square in that 2-row can either be in column 3 or column 4. Once we determine that, the other shaded squares are uniquely determined. Thus, there are ${3 \choose 2}{3 \choose 1}\times 2=18$ valid shadings in this case.
                \end{itemize}


                
                \begin{itemize}
                    \item If 1 shaded square in the first column is in a 1-row, and the other 2 are in 2-rows, then the 2-row without a shaded square in the first column must have shaded squares in both column 2 and column 3. This leaves 4 possible squares in the second column to be shaded (since there can't be another shaded square in the occupied 1-row). Thus, there are ${3 \choose 1}{3 \choose 2}\times {4 \choose 2}=54$ valid shadings in this case. (We only need to choose 2, since 1 of the shaded squares in the third column must go to the unoccupied 2-row).
                \end{itemize}


                
                \begin{itemize}
                    \item If all of the shaded squares in the first column are in the 3 2-rows, then if we choose any 3 squares in the second column to be shaded, then the third column is uniquely determined to create a valid shading. Thus, there are ${3 \choose 0}{3 \choose 3}\times{6 \choose 3} = 20$ valid shadings in this case.
                \end{itemize}


                In total, we have ${6\choose3}(1+18+54+20)=20*93=1860$. Thus our answer is $\boxed{860}$.

 





\textbf{Solution 8}

We can use generating functions. Suppose that the variables $a$, $b$, $c$, and $d$ represent shading a square that appears in the first, second, third, or fourth columns, respectively. Then if two squares in the row are shaded, then the row is represented by the generating function $ab+ac+ad+bc+bd+cd$, which we can write as $P(a,b,c,d)=(ab+cd)+(a+b)(c+d)$. Therefore, $P(a,b,c,d)^6$ represents all of the possible ways to color six rows such that each row has two shaded squares. We only want the possibilities when each column has three shaded squares, or rather, the coefficient of $a^3b^3c^3d^3$ in $P(a,b,c,d)^6$.

 By the Binomial Theorem,
 \[P(a,b,c,d)^6=\sum_{k=0}^6 \binom{6}{k} (ab+cd)^k(a+b)^{6-k}(c+d)^{6-k}.\tag{1}\]If we expand $(ab+cd)^k$, then the powers of $a$ and $b$ are always equal. Therefore, to obtain terms of the form $a^3b^3c^3d^3$, the powers of $a$ and $b$ in $(a+b)^{6-k}$ must be equal. In particular, only the central term in the binomial expansion will contribute, and this implies that $k$ must be even. We can use the same logic for $c$ and $d$. Therefore, the coefficient of $a^3b^3c^3d^3$ in the following expression is the same as the coefficient of $a^3b^3c^3d^3$ in (1).
 \[\sum_{k=0}^3 \binom{6}{2k} (ab+cd)^{2k}(ab)^{3-k}(cd)^{3-k}\binom{6-2k}{3-k}^2.\tag{2}\]Now we notice that the only way to obtain terms of the form $a^3b^3c^3d^3$ is if we take the central term in the binomial expansion of $(ab+cd)^{2k}$. Therefore, the terms that contribute to the coefficient of $a^3b^3c^3d^3$ in (2) are
 \[\sum_{k=0}^3 \binom{6}{2k}\binom{2k}{k}\binom{6-2k}{3-k}^2(abcd)^3.\]This sum is $400+1080+360+20=1860$ so the answer is $\boxed{860}$.

 




	\item (\textit{AIME I 2007}) For each positive integer $p$, let $b(p)$ denote the unique positive integer $k$ such that $|k-\sqrt{p}| < \frac{1}{2}$.  For example, $b(6) = 2$ and $b(23) = 5$.  If $S = \sum_{p=1}^{2007} b(p),$ find the remainder when $S$ is divided by 1000.


 
 \textbf{Answer: }955



\textbf{Solution 1}

$\left(k- \frac 12\right)^2=k^2-k+\frac 14$ and $\left(k+ \frac 12\right)^2=k^2+k+ \frac 14$. Therefore $b(p)=k$ if and only if $p$ is in this range, or $k^2-k<p\leq k^2+k$. There are $2k$ numbers in this range, so the sum of $b(p)$ over this range is $(2k)k=2k^2$. $44<\sqrt{2007}<45$, so all numbers $1$ to $44$ have their full range. Summing this up with the formula for the sum of the first $n$ squares ($\frac{n(n+1)(2n+1)}{6}$), we get $\sum_{k=1}^{44}2k^2=2\frac{44(44+1)(2*44+1)}{6}=58740$. We need only consider the $740$ because we are working with modulo $1000$.

 Now consider the range of numbers such that $b(p)=45$. These numbers are $\left\lceil\frac{44^2 + 45^2}{2}\right\rceil = 1981$ to $2007$. There are $2007 - 1981 + 1 = 27$ (1 to be inclusive) of them. $27*45=1215$, and $215+740= \boxed{955}$, the answer.

 





\textbf{Solution 2}

Let $p$ be in the range of $a^2 \le p < (a+1)^2$. Then, we need to find the point where the value of $b(p)$ flips from $k$ to $k+1$. This will happen when $p$ exceeds $(a+\frac{1}{2})^2$ or $a(a+1)+\frac{1}{4}$. Thus, if $a^2 \le p \le a(a+1)$ then $b(p)=a$. For $a(a+1) < p < (a+1)^2$, then $b(p)=a+1$. There are $a+1$ terms in the first set of $p$, and $a$ terms in the second set. Thus, the sum of $b(p)$ from $a^2 \le p <(a+1)^2$ is $2a(a+1)$ or $4\cdot\binom{a+1}{2}$. For the time being, consider that $S = \sum_{p=1}^{44^2-1} b(p)$. Then, the sum of the values of $b(p)$ is $4\binom{2}{2}+4\binom{3}{2}+\cdots +4\binom{44}{2}=4\left(\binom{2}{2}+\binom{3}{2}+\cdots +\binom{44}{2}\right)$. We can collapse this to $4\binom{45}{3}=56760$. Now, we have to consider $p$ from $44^2 \le p < 2007$. Considering $p$ from just $44^2 \le p \le 1980$, we see that all of these values have $b(p)=44$. Because there are $45$ values of $p$ in that range, the sum of $b(p)$ in that range is $45\cdot44=1980$. Adding this to $56760$ we get $58740$ or $740$ mod $1000$. Now, take the range $1980 < p \le 2007$. There are $27$ values of $p$ in this range, and each has $b(p)=45$. Thus, that contributes $27*45=1215$ or $215$ to the sum. Finally, adding $740$ and $215$ we get $740+215=\boxed{955}$.

 


 \href{https://mathproblemsolvingskills.wordpress.com/}{MathProblemSolvingSkills.com}

 


 




	\item (\textit{AIME I 2007}) In isosceles triangle $ABC$, $A$ is located at the origin and $B$ is located at (20,0).  Point $C$ is in the first quadrant with $AC = BC$ and angle $BAC = 75^{\circ}$.  If triangle $ABC$ is rotated counterclockwise about point $A$ until the image of $C$ lies on the positive $y$-axis, the area of the region common to the original and the rotated triangle is in the form $p\sqrt{2} + q\sqrt{3} + r\sqrt{6} + s$, where $p,q,r,s$ are integers.  Find $(p-q+r-s)/2$.


 
 \textbf{Answer: }875



\textbf{Solution}


                    \begin{figure}[H]
                        \centering
                        \includegraphics[width=0.5\textwidth]{Solution/AIME_2007_I_12_Solution_Figure_1.png}
                    \end{figure}
                    

 


 





\textbf{Solution 1}

Let the new triangle be $\triangle AB'C'$ ($A$, the origin, is a vertex of both triangles). Let $\overline{B'C'}$ intersect with $\overline{AC}$ at point $D$, $\overline{BC}$ intersect with $\overline{B'C'}$ at $E$, and $\overline{BC}$ intersect with $\overline{AB'}$ at $F$. The region common to both triangles is the \href{https://artofproblemsolving.com/wiki/index.php?title=Quadrilateral}{quadrilateral} $ADEF$. Notice that $[ADEF] = [\triangle ADB'] - [\triangle EFB']$, where we let $[\ldots]$ denote area. 

 To find $[\triangle ADB']$:Since $\angle B'AC'$ and $\angle BAC$ both have measures $75^{\circ}$, both of their \href{https://artofproblemsolving.com/wiki/index.php?title=Complement}{complements} are $15^{\circ}$, and $\angle DAB' = 90 - 2(15) = 60^{\circ}$. We know that $\angle DB'A = 75^{\circ}$, so $\angle ADB' = 180 - 60 - 75 = 45^{\circ}$.

 Thus $\triangle ADB'$ is a $45 - 60 - 75 \triangle$. It can be solved by drawing an \href{https://artofproblemsolving.com/wiki/index.php?title=Altitude}{altitude} splitting the $75^{\circ}$ angle into $30^{\circ}$ and $45^{\circ}$ angles, forming a $30-60-90$ \href{https://artofproblemsolving.com/wiki/index.php?title=Right\_triangle}{right triangle} and a $45-45-90$ isosceles right triangle. Since we know that $AB' = 20$, the base of the $30-60-90$ triangle is $10$, the base of the $45-45-90$ is $10\sqrt{3}$, and their common height is $10\sqrt{3}$. Thus, the total area of $[\triangle ADB'] = \frac{1}{2}(10\sqrt{3})(10\sqrt{3} + 10) = \boxed{150 + 50\sqrt{3}}$.

 To find $[\triangle EFB']$:Since $\triangle AFB$ is also a $15-75-90$ triangle,

 $AF = 20\sin 75 = 20 \sin (30 + 45) = 20\left(\frac{\sqrt{2} + \sqrt{6}}4\right) = 5\sqrt{2} + 5\sqrt{6}$and

 $FB' = AB' - AF = 20 - 5\sqrt{2} - 5\sqrt{6}$Since $[\triangle EFB'] = \frac{1}{2} (FB' \cdot EF) = \frac{1}{2} (FB') (FB' \tan 75^{\circ})$. With some horrendous \href{https://artofproblemsolving.com/wiki/index.php?title=Algebra}{algebra}, we can calculate
 \begin{align*} [\triangle EFB'] &= \frac{1}{2}\tan (30 + 45) \cdot (20 - 5\sqrt{2} - 5\sqrt{6})^2 \\ &= 25 \left(\frac{\frac{1}{\sqrt{3}} + 1}{1 - \frac{1}{\sqrt{3}}}\right) \left(8 - 2\sqrt{2} - 2\sqrt{6} - 2\sqrt{2} + 1 + \sqrt{3} - 2\sqrt{6} + \sqrt{3} + 3\right) \\ &= 25(2 + \sqrt{3})(12 - 4\sqrt{2} - 4\sqrt{6} + 2\sqrt{3}) \\ [\triangle EFB'] &= \boxed{- 500\sqrt{2} + 400\sqrt{3} - 300\sqrt{6} +750}. \end{align*}


 To finish,
 \begin{align*} [ADEF] &= [\triangle ADB'] - [\triangle EFB']\\        &= \left(150 + 50\sqrt{3}\right) - \left(-500\sqrt{2} + 400\sqrt{3} - 300\sqrt{6} + 750\right)\\        &=500\sqrt{2} - 350\sqrt{3} + 300\sqrt{6} - 600\\ \end{align*}
Hence, $\frac{p-q+r-s}{2} = \frac{500 + 350 + 300 + 600}2 = \frac{1750}2 = \boxed{\boxed{875}}$.

 





\textbf{Solution 2}

Redefine the points in the same manner as the last time ($\triangle AB'C'$, intersect at $D$, $E$, and $F$). This time, notice that $[ADEF] = [\triangle AB'C'] - ([\triangle ADC'] + [\triangle EFB'])$.

 The area of $[\triangle AB'C'] = [\triangle ABC]$. The \href{https://artofproblemsolving.com/wiki/index.php?title=Altitude}{altitude} of $\triangle ABC$ is clearly $10 \tan 75 = 10 \tan (30 + 45)$. The \href{https://artofproblemsolving.com/wiki/index.php/Trigonometric\_identities\#Angle\_addition\_identities}{tangent addition rule} yields $10(2 + \sqrt{3})$ (see above). Thus, $[\triangle ABC] = \frac{1}{2} 20 \cdot (20 + 10\sqrt{3}) = 200 + 100\sqrt{3}$. 

 The area of $[\triangle ADC']$ (with a side on the y-axis) can be found by splitting it into two triangles, $30-60-90$ and $15-75-90$ \href{https://artofproblemsolving.com/wiki/index.php?title=Right\_triangle}{right triangles}. $AC' = AC = \frac{10}{\sin 15}$. The \href{https://artofproblemsolving.com/wiki/index.php?title=Trigonometric\_identity}{sine subtraction rule} shows that $\frac{10}{\sin 15} = \frac{10}{\frac{\sqrt{6} - \sqrt{2}}4} = \frac{40}{\sqrt{6} - \sqrt{2}} = 10(\sqrt{6} + \sqrt{2})$. $AC'$, in terms of the height of $\triangle ADC'$, is equal to $h(\sqrt{3} + \tan 75) = h(\sqrt{3} + 2 + \sqrt{3})$. 
 \begin{align*} [ADC'] &= \frac 12 AC' \cdot h \\ &= \frac 12 (10\sqrt{6} + 10\sqrt{2})\left(\frac{10\sqrt{6} + 10\sqrt{2}}{2\sqrt{3} + 2}\right) \\ &= \frac{(800 + 400\sqrt{3})}{(2 + \sqrt{3})}\cdot\frac{2 - \sqrt{3}}{2-\sqrt{3}} \\ &= \frac{400\sqrt{3} + 400}8 = 50\sqrt{3} + 50 \end{align*}


 The area of $[\triangle EFB']$ was found in the previous solution to be $- 500\sqrt{2} + 400\sqrt{3} - 300\sqrt{6} +750$. 

 Therefore, $[ADEF]$ $= (200 + 100\sqrt{3}) - \left((50 + 50\sqrt{3}) + (-500\sqrt{2} + 400\sqrt{3} - 300\sqrt{6} +750)\right)$ $= 500\sqrt{2} - 350\sqrt{3} + 300\sqrt{6} - 600$, and our answer is $\boxed{875}$.

 





\textbf{Solution 3}


                    \begin{figure}[H]
                        \centering
                        \includegraphics[width=0.5\textwidth]{Solution/AIME_2007_I_12_Solution_Figure_1.png}
                    \end{figure}
                    Call the points of the intersections of the triangles $D$, $E$, and $F$ as noted in the diagram (the points are different from those in the diagram for solution 1). $\overline{AD}$ \href{https://artofproblemsolving.com/wiki/index.php?title=Bisect}{bisects} $\angle EDE'$.

 Through \href{https://artofproblemsolving.com/wiki/index.php/Congruent\_(geometry)\#HL\_Congruence}{HL} congruency, we can find that $\triangle AED$ is \href{https://artofproblemsolving.com/wiki/index.php?title=Congruent}{congruent} to $\triangle AE'D$. This divides the region $AEDF$ (which we are trying to solve for) into two congruent triangles and an \href{https://artofproblemsolving.com/wiki/index.php?title=Isosceles\_triangle}{isosceles} \href{https://artofproblemsolving.com/wiki/index.php?title=Right\_triangle}{right triangle}.  

 $AE = 20 \cos 15 = 20 \cos (45 - 30) = 20 \cdot \frac{\sqrt{6} + \sqrt{2}}{4} = 5\sqrt{6} + 5\sqrt{2}$Since $FE' = AE' = AE$, we find that $[AE'F] = \frac 12 (5\sqrt{6} + 5\sqrt{2})^2 = 100 + 50\sqrt{3}$.

 Now, we need to find $[AED] = [AE'D]$. The acute angles of the triangles are $\frac{15}{2}$ and $90 - \frac{15}{2}$. By repeated application of the \href{https://artofproblemsolving.com/wiki/index.php?title=Trigonometric\_identity}{half-angle formula}, we can find that $\tan \frac{15}{2} = \sqrt{2} - \sqrt{3} + \sqrt{6} - 2$.

 The area of $[AED] = \frac 12 \left(20 \cos 15\right)^2 \left(\tan \frac{15}{2}\right)$. Thus, $[AED] + [AE'D] = 2\left(\frac 12((5\sqrt{6} + 5\sqrt{2})^2 \cdot (\sqrt{2} - \sqrt{3} + \sqrt{6} - 2))\right)$, which eventually simplifies to $500\sqrt{2} - 350\sqrt{3} + 300\sqrt{6} - 600$. 

 Adding them together, we find that the solution is $[AEDF] = [AE'F] + [AED] + [AE'D]$ $= 100 + 50\sqrt{3} + 500\sqrt{2} - 350\sqrt{3} + 300\sqrt{6} - 600=$ $= 500\sqrt{2} - 350\sqrt{3} + 300\sqrt{6} - 600$, and the answer is $\boxed{875}$.

 





\textbf{Solution 4}

From the given information, calculate the coordinates of all the points (by finding the equations of the lines, equating). Then, use the \href{https://artofproblemsolving.com/wiki/index.php?title=Shoelace}{shoelace} method to calculate the area of the intersection.

 
                \begin{itemize}
                    \item $\overline{AC}$: $y = (\tan 75) x = (2 + \sqrt{3})x$\item $\overline{AB'}$: $y = (\tan 15) x = (2 - \sqrt{3})x$\item $\overline{BC}$: It passes thru $(20,0)$, and has a slope of $-\tan75 = -(2 + \sqrt{3})$. The equation of its line is  $y = (2+\sqrt{3})(20 - x)$.\item $\overline{B'C'}$: $AC' = AC = \frac{10}{\cos 75} = 10\sqrt{6} + 10\sqrt{2}$, so it passes thru point $(0, 10\sqrt{6} + 10\sqrt{2})$. It has a slope of $-\tan 60 = -\sqrt{3}$. So the equation of its line is $y = -\sqrt{3}x + 10(\sqrt{6} + \sqrt{2})$.
                \end{itemize}


                Now, we can equate the equations to find the intersections of all the points.

 
                \begin{itemize}
                    \item $A (0,0)$
                \end{itemize}


                
                \begin{itemize}
                    \item $D$ is the intersection of $\overline{BC},\ \overline{B'C'}$. $(2+\sqrt{3})(20-x) = -\sqrt{3}x + 10(\sqrt{6} + \sqrt{2})$. Therefore, $x = 5(4 + 2\sqrt{3}-\sqrt{6}-\sqrt{2})$, $y = 5(3\sqrt{6} + 5\sqrt{2}-4\sqrt{3}-6)$.
                \end{itemize}


                
                \begin{itemize}
                    \item $E$ is the intersection of $\overline{AB'},\ \overline{BC}$. $(2 - \sqrt{3})x =(2+\sqrt{3})(20-x)$. Therefore, $x = 5(2+\sqrt{3})$, $y = 5$.
                \end{itemize}


                
                \begin{itemize}
                    \item $F$ is the intersection of $\overline{AC},\ \overline{B'C'}$. $(2+\sqrt{3})x = -\sqrt{3}x + 10(\sqrt{6} + \sqrt{2})$. Therefore, $x = 5\sqrt{2}$, $y = 10\sqrt{2}+ 5\sqrt{6}$.
                \end{itemize}


                We take these points and tie them together by shoelace, and the answer should come out to be $\boxed{875}$.

 


 \href{https://mathproblemsolvingskills.wordpress.com/}{MathProblemSolvingSkills.com}

 


 




	\item (\textit{AIME I 2007}) A square pyramid with base $ABCD$ and vertex $E$ has eight edges of length 4.  A plane passes through the midpoints of $AE$, $BC$, and $CD$.  The plane's intersection with the pyramid has an area that can be expressed as $\sqrt{p}$.  Find $p$.


 
 \textbf{Answer: }080



\textbf{Solution 1}

Note first that the intersection is a \href{https://artofproblemsolving.com/wiki/index.php?title=Pentagon}{pentagon}.

 Use 3D analytical geometry, setting the origin as the center of the square base and the pyramid’s points oriented as shown above. $A(-2,2,0),\ B(2,2,0),\ C(2,-2,0),\ D(-2,-2,0),\ E(0,0,2\sqrt{2})$. Using the coordinates of the three points of intersection $(-1,1,\sqrt{2}),\ (2,0,0),\ (0,-2,0)$, it is possible to determine the equation of the plane. The equation of a plane resembles $ax + by + cz = d$, and using the points we find that $2a = d \Longrightarrow d = \frac{a}{2}$, $-2b = d \Longrightarrow d = \frac{-b}{2}$, and $-a + b + \sqrt{2}c = d \Longrightarrow -\frac{d}{2} - \frac{d}{2} + \sqrt{2}c = d \Longrightarrow c = d\sqrt{2}$. It is then $x - y + 2\sqrt{2}z = 2$.

 <center> 
                    \begin{figure}[H]
                        \centering
                        \includegraphics[width=0.5\textwidth]{Solution/AIME_2007_I_13_Solution_Figure_1.png}
                    \end{figure}
                    
                    \begin{figure}[H]
                        \centering
                        \includegraphics[width=0.5\textwidth]{Solution/AIME_2007_I_13_Solution_Figure_1.png}
                    \end{figure}
                    

 </center>Write the equation of the lines and substitute to find that the other two points of intersection on $\overline{BE}$, $\overline{DE}$ are $\left(\frac{\pm 3}{2},\frac{\pm 3}{2},\frac{\sqrt{2}}{2}\right)$. To find the area of the pentagon, break it up into pieces (an \href{https://artofproblemsolving.com/wiki/index.php?title=Isosceles\_triangle}{isosceles triangle} on the top, an \href{https://artofproblemsolving.com/wiki/index.php?title=Isosceles\_trapezoid}{isosceles trapezoid} on the bottom). Using the \href{https://artofproblemsolving.com/wiki/index.php?title=Distance\_formula}{distance formula} ($\sqrt{a^2 + b^2 + c^2}$), it is possible to find that the area of the triangle is $\frac{1}{2}bh \Longrightarrow \frac{1}{2} 3\sqrt{2} \cdot \sqrt{\frac 52} = \frac{3\sqrt{5}}{2}$. The trapezoid has area $\frac{1}{2}h(b_1 + b_2) \Longrightarrow \frac 12\sqrt{\frac 52}\left(2\sqrt{2} + 3\sqrt{2}\right) = \frac{5\sqrt{5}}{2}$. In total, the area is $4\sqrt{5} = \sqrt{80}$, and the solution is $\boxed{080}$.

 





\textbf{Solution 2}

Use the same coordinate system as above, and let the plane determined by $\triangle PQR$ intersect $\overline{BE}$ at $X$ and $\overline{DE}$ at $Y$. Then the line $\overline{XY}$ is the intersection of the planes determined by $\triangle PQR$ and $\triangle BDE$.

 Note that the plane determined by $\triangle BDE$ has the equation $x=y$, and $\overline{PQ}$ can be described by $x=2(1-t)-t,\ y=t,\ z=t\sqrt{2}$. It intersects the plane when $2(1-t)-t=t$, or $t=\frac{1}{2}$. This intersection point has $z=\frac{\sqrt{2}}{2}$. Similarly, the intersection between $\overline{PR}$ and $\triangle BDE$ has $z=\frac{\sqrt{2}}{2}$. So $\overline{XY}$ lies on the plane $z=\frac{\sqrt{2}}{2}$, from which we obtain $X=\left( \frac{3}{2},\frac{3}{2},\frac{\sqrt{2}}{2}\right)$ and $Y=\left( -\frac{3}{2},-\frac{3}{2},\frac{\sqrt{2}}{2}\right)$. The area of the pentagon $EXQRY$ can be computed in the same way as above.

 





\textbf{Solution 3}

<center> 
                    \begin{figure}[H]
                        \centering
                        \includegraphics[width=0.5\textwidth]{Solution/AIME_2007_I_13_Solution_Figure_1.png}
                    \end{figure}
                    

 </center>Extend $\overline{RQ}$ and $\overline{AB}$. The point of intersection is $H$. Connect $\overline{PH}$. $\overline{EB}$ intersects $\overline{PH}$ at $X$. Do the same for $\overline{QR}$ and $\overline{AD}$, and let the intersections be $I$ and $Y$

 Because $Q$ is the midpoint of $\overline{BC}$, and $\overline{AB}\parallel\overline{DC}$, so $\triangle{RQC}\cong\triangle{HQB}$. $\overline{BH}=2$. 

 Because $\overline{BH}=2$, we can use mass point geometry to get that $\overline{PX}=\overline{XH}$. $|\triangle{XHQ}|=\frac{\overline{XH}}{\overline{PH}}\cdot\frac{\overline{QH}}{\overline{HI}}\cdot|\triangle{PHI}|=\frac{1}{6}\cdot|\triangle{PHI}|$ 

 Using the same principle, we can get that $|\triangle{IYR}|=\frac{1}{6}|\triangle{PHI}|$

 Therefore, the area of $PYRQX$ is $\frac{2}{3}\cdot|\triangle{PHI}|$

 $\overline{RQ}=2\sqrt{2}$, so $\overline{IH}=6\sqrt{2}$. Using the law of cosines, $\overline{PH}=\sqrt{28}$. The area of $\triangle{PHI}=\frac{1}{2}\cdot\sqrt{28-18}\cdot6\sqrt{2}=6\sqrt{5}$

 Using this, we can get the area of $PYRQX = \sqrt{80}$ so the answer is $\fbox{080}$.

 


 \href{https://mathproblemsolvingskills.wordpress.com/}{MathProblemSolvingSkills.com}

 


 




	\item (\textit{AIME I 2007}) A sequence is defined over non-negative integral indexes in the following way: $a_{0}=a_{1}=3$, $a_{n+1}a_{n-1}=a_{n}^{2}+2007$.


 Find the greatest integer that does not exceed $\frac{a_{2006}^{2}+a_{2007}^{2}}{a_{2006}a_{2007}}$


 
 \textbf{Answer: }224



\textbf{Solution 1}

Define a function $f(n)$ on the non-negative integers, as 
 \[f(n) = \frac{a_n^2 + a_{n+1}^2}{a_na_{n+1}} = \frac{a_n}{a_{n+1}}+\frac{a_{n+1}}{a_{n}}\]We want $\left\lfloor f(2006) \right\rfloor$. 

 Consider the relation $a_{n+1}a_{n-1}=a_{n}^{2}+2007$. Dividing through by $a_{n}a_{n-1}$, we get 
 \[.\phantom{------------} \frac{a_{n+1}}{a_{n}} = \frac{a_{n}}{a_{n-1}} + \frac{2007}{a_{n}a_{n-1}} \phantom{------------} (1)\]and multiplying both sides by $\frac{a_{n-1}}{a_{n+1}}$, we get 
 \[.\phantom{------------} \frac{a_{n-1}}{a_{n}} = \frac{a_{n}}{a_{n+1}} + \frac{2007}{a_{n}a_{n+1}} \phantom{------------} (2)\]Adding LHS of $(1)$ with RHS of $(2)$ (and vice-versa), we get
 \[\frac{a_{n+1}}{a_{n}} + \frac{a_{n}}{a_{n+1}} + \frac{2007}{a_{n}a_{n+1}} = \frac{a_{n}}{a_{n-1}} + \frac{a_{n-1}}{a_{n}} + \frac{2007}{a_{n}a_{n-1}}\]i.e.
 \[f(n)+ \frac{2007}{a_{n}a_{n+1}} = f(n-1) + \frac{2007}{a_{n}a_{n-1}}\]Rearranging this to the form $f(n)-f(n-1)$ and summing over $n=1$ to $n=2006$, we notice that most of the terms on each side telescope. Without rearranging, you can see that most terms cancel against the corresponding term on the other side. We are left with
 \[f(2006) + \frac{2007}{a_{2006}a_{2007}} = f(0) + \frac{2007}{a_{1}a_{0}}\]We have $f(0) = 2$, and $2007/a_0a_1 = 2007/9 = 223$. So
 \[f(2006) = 2 + 223 - \frac{2007}{a_{2006}a_{2007}} = 224 + \left( 1 - \frac{2007}{a_{2006}a_{2007}}\right)\]Since all the $a_i$ are positive, $(1)$ tells us that the ratio $a_{n+1}/a_n$ of successive terms is increasing. Since this ratio starts with $a_1/a_0 = 1$, this means that the sequence $(a_n)$ is increasing. Since $a_2=672$ already, we must have $a_{2006}a_{2007} > 672^2 > 2007$. It follows that $0<1-\frac{2007}{a_{2006}a_{2007}}<1$ and so $\left\lfloor f(2006) \right\rfloor = \boxed{224}$.

 





\textbf{Solution 2}

We are given that

 $a_{n+1}a_{n-1}= a_{n}^{2}+2007$,$a_{n-1}^{2}+2007 = a_{n}a_{n-2}$.

 Add these two equations to get

 
            \begin{enumerate}[label=(\roman*)]
                \item $a_{n-1}(a_{n-1}+a_{n+1}) = a_{n}(a_{n}+a_{n-2})$
            \end{enumerate}
            
            \begin{enumerate}[label=(\roman*)]
                \item $\frac{a_{n+1}+a_{n-1}}{a_{n}}= \frac{a_{n}+a_{n-2}}{a_{n-1}}$.
            \end{enumerate}
            This is an \href{https://artofproblemsolving.com/wiki/index.php?title=Invariant}{invariant}. Defining $b_{i}= \frac{a_{i}}{a_{i-1}}$ for each $i \ge 2$, the above equation means

 $b_{n+1}+\frac{1}{b_{n}}= b_{n}+\frac{1}{b_{n-1}}$.

 We can thus calculate that $b_{2007}+\frac{1}{b_{2006}}= b_{3}+\frac{1}{b_{2}}= 225$. Using the equation $a_{2007}a_{2005}=a_{2006}^{2}+2007$ and dividing both sides by $a_{2006}a_{2005}$, notice that $b_{2007}= \frac{a_{2007}}{a_{2006}}= \frac{a_{2006}^{2}+2007}{a_{2006}a_{2005}}> \frac{a_{2006}}{a_{2005}}= b_{2006}$. This means that

 $b_{2007}+\frac{1}{b_{2007}}< b_{2007}+\frac{1}{b_{2006}}= 225$. It is only a tiny bit less because all the $b_i$ are greater than $1$, so we conclude that the floor of $\frac{a_{2007}^{2}+a_{2006}^{2}}{a_{2007}a_{2006}}= b_{2007}+\frac{1}{b_{2007}}$ is $\boxed{224}$.

 





\textbf{Solution 3}

The equation $a_{n+1}a_{n-1}-a_n^2=2007$ looks like the determinant 
 \[\left|\begin{array}{cc}a_{n+1}&a_n\\a_n&a_{n-1}\end{array}\right|=2007.\]Therefore, the determinant of this matrix is invariant. Guessing that this sequence might be a linear recursion because of the matrix form given below, we define the sequence $b_n$ defined by $b_1=b_2=3$ and $b_{n+1}=\alpha b_n+\beta b_{n-1}$ for $n\ge 2$. We wish to find $\alpha$ and $\beta$ such that $a_n=b_n$ for all $n\ge 1$. To do this, we use the following matrix form of a linear recurrence relation

 
 \[\left(\begin{array}{cc}b_{n+1}&b_n\\b_n&b_{n-1}\end{array}\right)=\left(\begin{array}{cc}\alpha&\beta\\1&0\end{array}\right)\left(\begin{array}{cc}b_{n}&b_{n-1}\\b_{n-1}&b_{n-2}\end{array}\right).\] 

 When we take determinants, this equation becomes

 
 \[\text{det}\left(\begin{array}{cc}b_{n+1}&b_n\\b_n&b_{n-1}\end{array}\right)=\text{det}\left(\begin{array}{cc}\alpha&\beta\\1&0\end{array}\right)\text{det}\left(\begin{array}{cc}b_{n}&b_{n-1}\\b_{n-1}&b_{n-2}\end{array}\right).\]

 We want 
 \[\text{det}\left(\begin{array}{cc}b_{n+1}&b_n\\b_n&b_{n-1}\end{array}\right)=2007\] for all $n$. Therefore, we replace the two matrices by $2007$ to find that 

 
 \[2007=\text{det}\left(\begin{array}{cc}\alpha&\beta\\1&0\end{array}\right)\cdot 2007\]
 \[1=\text{det}\left(\begin{array}{cc}\alpha&\beta\\1&0\end{array}\right)=-\beta.\]Therefore, $\beta=-1$. Computing that $a_3=672$, and using the fact that $b_3=\alpha b_2-b_1$, we conclude that $\alpha=225$. Clearly, $a_1=b_1$, $a_2=b_2$, and $a_3=b_3$. We claim that $a_n=b_n$ for all $n\ge 1$. We proceed by \href{https://artofproblemsolving.com/wiki/index.php?title=Induction}{induction}. If $a_k=b_k$ for all $k\le n$, then clearly, 
 \[b_nb_{n-2}-b_{n-1}^2=a_na_{n-2}-a_{n-1}^2=2007.\] We also know by the definition of $b_{n+1}$ that

 
 \[\text{det}\left(\begin{array}{cc}b_{n+1}&b_n\\b_n&b_{n-1}\end{array}\right)=\text{det}\left(\begin{array}{cc}\alpha&\beta\\1&0\end{array}\right)\text{det}\left(\begin{array}{cc}b_{n}&b_{n-1}\\b_{n-1}&b_{n-2}\end{array}\right).\]

 We know that the RHS is $2007$ by previous work. Therefore, $b_{n+1}b_{n-1}-b_n^2=2007$. After substuting in the values we know, this becomes $b_{n+1}a_{n-1}-a_n^2=2007$.  Thinking of this as a linear equation in the variable $b_{n+1}$, we already know that this has the solution $b_{n+1}=a_{n+1}$. Therefore, by induction, $a_n=b_n$ for all $n\ge 1$. We conclude that $a_n$ satisfies the linear recurrence $a_{n+1}=225a_n-a_{n-1}$.

 It's easy to prove that $a_n$ is a strictly increasing sequence of integers for $n\ge 3$. Now 

 
 \[\frac{a_{2007}^2+a_{2006}^2}{a_{2007}a_{2006}}=\frac{a_{2007}}{a_{2006}}+\frac{a_{2006}}{a_{2007}}=\frac{225a_{2006}-a_{2005}}{a_{2006}}+\frac{a_{2006}}{a_{2007}}.\]

 
 \[=225+\frac{a_{2006}}{a_{2007}}-\frac{a_{2005}}{a_{2006}}=225+\frac{a_{2006}^2-a_{2005}a_{2007}}{a_{2005}a_{2006}}.\]

 
 \[=225-\frac{2007}{a_{2005}a_{2006}}.\]

 The sequence certainly grows fast enough such that $\frac{2007}{a_{2005}a_{2006}}<1$. Therefore, the largest integer less than or equal to this value is $\boxed{224}$.

 





\textbf{Solution 4 ( generalized )}

This is a more elementary and rigorous solution to a slightly generalized version. The defining recursive sequence is generalized to 

 
 \[a_{n+1}a_{n-1} = a_n^2 + 9k, ---------(1)\]where $k$ is a positive integer and $a_0 = a_1 = 3.$

 Lemma 1  : For $n \geq 1$,
 \[a_{n+1} = ( k + 2)a_n - a_{n-1}. ---------(2)\]We shall prove by induction. From (1), $a_2 = 3k + 3$. From the lemma, $a_2 = (k + 2) 3 - 3 = 3k + 3.$ Base case proven. Assume that the lemma is true for some $t \geq 1$. Then, eliminating the $a_{t-1}$ using (1) and (2) gives

 
 \[(k+2)a_ta_{t+1} = a_t^2 + a_{t+1}^2 + 9k. ---------(3)\]

 It follows from (2) that 

 
 \[(k+2)a_{t+1} - a_t  =\frac{(k+2)a_{t+1}a_t - a_t^2}{a_t}  =\frac{a_{t+1}^2 + 9k}{a_t} =a_{t+2},\]

 where the last line followed from (1) for case $n = t+1$.

 Lemma 2 : For $n \geq 0,$
 \[a_{n+1} \geq a_{n}.\]Base case is obvious. Assume that $a_{t+1} \geq a_{t}$ for some $t \geq 0$. Then it follows that 
 \[a_{t+2} =\frac{a_{t+1}^2 + 9k}{a_t}  = a_{t+1}(\frac{a_{t+1}}{a_t} ) + 9k \geq a_{t+1} + 9k  > a_{t+1}.\]

 This completes the induction.

 Lemma 3 : For $n \geq 1,$
 \[a_n a_{n+1} > 9k\]

 Using (1) and Lemma 2, for $n \geq 1,$ 
 \[a_{n+1}a_n \geq a_{n+1}a_{n-1}  = a_n^2 + 9k  > 9k\]

 Finally, using (3), for $n \geq 1,$
 \[\frac{a_n^2 + a_{n+1}^2}{a_n a_{n+1}} =\frac{(k+2)a_n a_{n+1} - 9k}{a_n a_{n+1}} = k+2 -\frac{9k}{a_n a_{n+1}}.\]Using lemma 3, the largest integer less than or equal to this value would be $k + 1$.

 





\textbf{Solution 5 (pure algebra)}

We will try to manipulate $\frac{a_0^2+a_1^2}{a_0a_1}$ to get $\frac{a_1^2+a_2^2}{a_1a_2}$. 

 $\frac{a_0^2+a_1^2}{a_0a_1} = \frac{a_0+\frac{a_1^2}{a_0}}{a_1} = \frac{a_0+\frac{a_1^2+2007}{a_0}-\frac{2007}{a_0}}{a_1}$Using the recurrence relation, $= \frac{a_0+a_2-\frac{2007}{a_0}}{a_1}  = \frac{a_0a_2+a_2^2-\frac{2007a_2}{a_0}}{a_1a_2}$Applying the relation to $a_0a_2$, $= \frac{a_1^2+2007+a_2^2-\frac{2007a_2}{a_0}}{a_1a_2} = \frac{a_1^2+a_2^2}{a_1a_2}+\frac{2007}{a_1a_2}-\frac{2007}{a_0a_1}$

 We can keep on using this method to get that $\frac{a_0^2+a_1^2}{a_0a_1} = \frac{a_{2006}^2+a_{2007}^2}{a_{2006}a_{2007}}+\frac{2007}{a_{2006}a_{2007}}-\frac{2007}{a_{2005}a_{2006}}+\frac{2007}{a_{2005}a_{2006}}-\ldots-\frac{2007}{a_{0}a_{1}}$

 This telescopes to $\frac{a_0^2+a_1^2}{a_0a_1} = \frac{a_{2006}^2+a_{2007}^2}{a_{2006}a_{2007}}+\frac{2007}{a_{2006}a_{2007}}-\frac{2007}{a_{0}a_{1}}$

 or $\frac{a_{2006}^2+a_{2007}^2}{a_{2006}a_{2007}} = \frac{a_0^2+a_1^2}{a_0a_1}+\frac{2007}{a_{0}a_{1}}-\frac{2007}{a_{2006}a_{2007}}$

 Finding the first few values, we notice that they increase rapidly, so $\frac{2007}{a_{2006}a_{2007}} < 1$. Calculating the other values, $\frac{a_{2006}^2+a_{2007}^2}{a_{2006}a_{2007}} = 2+223-\frac{2007}{a_{2006}a_{2007}}$.

 The greatest number that does not exceed this is $\boxed{224}$

 





\textbf{Solution 6 (using limits)}

Let's start by computing the first couple terms of the given sequence so we know what we're working with. It's given that $a_0 = a_1 = 0$, and solving for $a_2$ and $a_3$ using the given relation we get $a_2 = 672 = 3(224)$ and $a_3 = 3(224^{2} + 223)$, respectively. It will be clear why I decided to factor these expressions as I did momentarily. 

 Next, let's see what the expression $\frac{a_{n}^{2}+a_{n + 1}^{2}}{a_{n}a_{n + 1}}$ looks like for small values of $n$. For $n = 1$, we get $\frac{1 + 224^2}{224}$, the floor of which is clearly $224$ because the $1$ in the numerator is insignificant. Repeating the procedure for $n + 1$ is somewhat messier, but we end up getting $\frac{224^4 + 224^2\cdot223\cdot2 + 224^2 + 223^2}{224^3 + 224\cdot223}$. It's not too hard to see that $224^4$ is much larger than the sum of the remaining terms in the numerator, and that $224^3$ is similarly a lot greater than the other term in the denominator. In fact, the second-largest term in the numerator is only barely larger than $224^3$, while the second-largest term in the denominator is smaller than $224^2$. Thus, the floor of this expression will come out to be $224$ as well.

 Now we must consider whether this finding holds true for the rest of the sequence we're examining. First of all, notice that each time $n$ increases by $1$, the degrees of both the numerator and denominator increase by $2$, because we are squaring the $n+1th$ term in the numerator while the same term, appearing in the denominator, is being generated in part by squaring the term before it (seen in the relation $a_{n+1}^2 = (\frac{a_{n}^2 + 2007}{a_{n-1}})^2$). Thus, the degree of the numerator will always be one greater than the degree of the denominator; we have only to show that the other terms in the expression remain sufficiently small enough so as not to disturb the $\approx224:1$ ratio between the two. 

 For the non-greatest terms in the expression to offset this ratio for values of $n$ in the ballpark of $2006$, they would have to have massive coefficients, because or else they are dwarfed by the additional $224$ attached to the leading term. Thankfully, this is not the case. Since we are simply squaring previous terms repeatedly to get new terms, we end up getting a sequence that is approximately equal to $\frac{224^k+224^{k-1}+224^{k-2} . . . }{224^{k-1}+224^{k-2} . . . }$ for all $k\geq2$, whose $\lim_{k\to\infty}=$ $\boxed{224}$.

 


 \href{https://mathproblemsolvingskills.wordpress.com/}{MathProblemSolvingSkills.com}

 


 




	\item (\textit{AIME I 2007}) Let $ABC$ be an equilateral triangle, and let $D$ and $F$ be points on sides $BC$ and $AB$, respectively, with $FA = 5$ and $CD = 2$.  Point $E$ lies on side $CA$ such that angle $DEF = 60^{\circ}$.  The area of triangle $DEF$ is $14\sqrt{3}$.  The two possible values of the length of side $AB$ are $p \pm q \sqrt{r}$, where $p$ and $q$ are rational, and $r$ is an integer not divisible by the square of a prime.  Find $r$.


 
 \textbf{Answer: }989



\textbf{Solution}


                    \begin{figure}[H]
                        \centering
                        \includegraphics[width=0.5\textwidth]{Solution/AIME_2007_I_15_Solution_Figure_1.png}
                    \end{figure}
                    

 Denote the length of a side of the triangle $x$, and of $\overline{AE}$ as $y$. The area of the entire equilateral triangle is $\frac{x^2\sqrt{3}}{4}$. Add up the areas of the triangles using the $\frac{1}{2}ab\sin C$ formula (notice that for the three outside triangles, $\sin 60 = \frac{\sqrt{3}}{2}$): $\frac{x^2\sqrt{3}}{4} = \frac{\sqrt{3}}{4}(5 \cdot y + (x - 2)(x - 5) + 2(x - y)) + 14\sqrt{3}$. This simplifies to $\frac{x^2\sqrt{3}}{4} = \frac{\sqrt{3}}{4}(5y + x^2 - 7x + 10 + 2x - 2y + 56)$. Some terms will cancel out, leaving $y = \frac{5}{3}x - 22$.

 $\angle FEC$ is an \href{https://artofproblemsolving.com/wiki/index.php?title=Exterior\_angle}{exterior angle} to $\triangle AEF$, from which we find that $60 + \angle CED = 60 + \angle AFE$, so $\angle CED = \angle AFE$. Similarly, we find that $\angle EDC = \angle AEF$. Thus, $\triangle AEF \sim \triangle CDE$. Setting up a \href{https://artofproblemsolving.com/wiki/index.php?title=Ratio}{ratio} of sides, we get that $\frac{5}{x-y} = \frac{y}{2}$. Using the previous relationship between $x$ and $y$, we can solve for $x$.

 $xy - y^2 = 10$

 $\frac{5}{3}x^2 - 22x - \left(\frac{5}{3}x - 22\right)^2 - 10 = 0$

 $\frac{5}{3}x^2 - \frac{25}{9}x^2 - 22x + 2 \cdot \frac{5 \cdot 22}{3}x - 22^2 - 10= 0$

 $10x^2 - 462x + 66^2 + 90 = 0$

 Use the \href{https://artofproblemsolving.com/wiki/index.php?title=Quadratic\_formula}{quadratic formula}, though we only need the root of the \href{https://artofproblemsolving.com/wiki/index.php?title=Discriminant}{discriminant}. This is $\sqrt{(7 \cdot 66)^2 - 4 \cdot 10 \cdot (66^2 + 90)} = \sqrt{49 \cdot 66^2 - 40 \cdot 66^2 - 4 \cdot 9 \cdot 100}$$= \sqrt{9 \cdot 4 \cdot 33^2 - 9 \cdot 4 \cdot 100} = 6\sqrt{33^2 - 100}$. The answer is $\boxed{989}$.

 





\textbf{Solution 2}

First of all, assume $EC=x,BD=m, ED=a, EF=b$, then we can find $BF=m-3, AE=2+m-x$It is not hard to find $ab*sin60^{\circ}*\frac{1}{2}=14\sqrt{3}, ab=56$, we apply LOC on $\triangle{DEF}, \triangle{BFD}$, getting that $(m-3)^2+m^2-m(m-3)=a^2+b^2-ab$, leads to $a^2+b^2=m^2-3m+65$Apply LOC on $\triangle{CED}, \triangle{AEF}$ separately, getting $4+x^2-2x=a^2; 25+(2+m-x)^2-5(2+m-x)=b^2.$ Add those terms together and use the equality $a^2+b^2=m^2-3m+65$, we can find:$2x^2-(2m+1)x+2m-42=0$

 According to basic angle chasing, $\angle{A}=\angle{C}; \angle{AFE}=\angle{CED}$, so $\triangle{AFE}\sim \triangle{CED}$, the ratio makes $\frac{5}{x}=\frac{2+m-x}{2}$, getting that $x^2-(2+m)x+10=0$Now we have two equations with $m$, and $x$ values for both equations must be the same, so we can solve for $x$ in two equations.$x=\frac{2m+1 \pm \sqrt{4m^2+4m+1-16m+336}}{4}; x=\frac{4+2m \pm \sqrt{4m^2+16m-144}}{4}$, then we can just use positive sign to solve, simplifies to $3+\sqrt{4m^2+16m-144}=\sqrt{4m^2-12m+337}$, getting $m=\frac{211-3\sqrt{989}}{10}$, since the triangle is equilateral, $AB=BC=2+m=\frac{231-3\sqrt{989}}{10}$, and the desired answer is $\boxed{989}$

 


 \href{https://mathproblemsolvingskills.wordpress.com/}{MathProblemSolvingSkills.com}

 


 




\end{enumerate}